% Fragment partage : inclus par liste-contacts.tex (dossier autonome)
% et par dossier-conception-alanya.tex (dossier client assemble).
% Ne pas y mettre de preambule ni de \part : c'est l'appelant qui les pose.

\chapter{Modèle de données}

\section{Deux tables}

Sur le modèle exact de \code{preferredContact} et \code{conv\_participants}
(\code{migrations/001\_initial\_schema.sql}).

\begin{lstlisting}[language=SQL, caption={\code{migrations/024\_contact\_lists.sql}}]
-- 024_contact_lists.sql -- listes de contacts + membres (plusieurs a plusieurs)

CREATE TABLE IF NOT EXISTS contact_list (
  idList     BIGINT       NOT NULL AUTO_INCREMENT,
  alanyaID   INT          NOT NULL,               -- proprietaire de la liste
  name       VARCHAR(60)  NOT NULL,
  color      VARCHAR(9)   NULL,                   -- #RRGGBB pour la puce
  created_at DATETIME     NOT NULL DEFAULT CURRENT_TIMESTAMP,
  PRIMARY KEY (idList),
  UNIQUE KEY uq_list_name (alanyaID, name),
  KEY idx_list_owner (alanyaID),
  CONSTRAINT fk_list_owner FOREIGN KEY (alanyaID)
    REFERENCES users(alanyaID) ON DELETE CASCADE
) ENGINE=InnoDB DEFAULT CHARSET=utf8mb4 COLLATE=utf8mb4_unicode_ci;

CREATE TABLE IF NOT EXISTS contact_list_member (
  idList     BIGINT   NOT NULL,
  idFriend   INT      NOT NULL,
  created_at DATETIME NOT NULL DEFAULT CURRENT_TIMESTAMP,
  PRIMARY KEY (idList, idFriend),
  KEY idx_clm_friend (idFriend),
  CONSTRAINT fk_clm_list   FOREIGN KEY (idList)
    REFERENCES contact_list(idList) ON DELETE CASCADE,
  CONSTRAINT fk_clm_friend FOREIGN KEY (idFriend)
    REFERENCES users(alanyaID) ON DELETE CASCADE
) ENGINE=InnoDB DEFAULT CHARSET=utf8mb4 COLLATE=utf8mb4_unicode_ci;
\end{lstlisting}

\begin{alerte}[Renumérotation --- 023 devient 024]
La spécification d'origine visait \code{023\_contact\_lists.sql}. Le numéro 023 est pris
par le socle de compte, qui doit être appliqué avant. Les migrations restent
\textbf{appliquées à la main}~: il n'y a ni exécuteur, ni table de suivi.
\end{alerte}

\begin{note}[Clé primaire composite plutôt qu'un identifiant technique]
\code{contact\_list\_member} n'a pas de colonne \code{id}~: sa clé est le couple
\code{(idList, idFriend)}. Cela interdit gratuitement le doublon d'appartenance --- il n'y
a rien à vérifier côté applicatif.
\end{note}

\section{L'appartenance n'est pas contrainte par la base}

La règle « un membre doit être un favori du propriétaire » ne peut pas s'exprimer en clé
étrangère~: elle porte sur le couple \code{(propriétaire de la liste, ami)}, pas sur
\code{idFriend} seul. Elle est donc vérifiée dans le contrôleur, à l'ajout, par une
requête sur \code{preferredContact}.

Conséquence à ne pas oublier~: si l'utilisateur \textbf{retire un favori}, ce contact reste
dans ses listes. Deux options --- le retirer en cascade applicative, ou le laisser et le
filtrer à la lecture. La première est plus propre et se code en trois lignes dans
\code{removePreferredContact}
(\code{src/controllers/preferredContactController.js}).

\chapter{API}

\section{Contrôleur --- \code{src/controllers/contactListController.js}}

Style copié de \code{preferredContactController.js}~: cadrage sur \code{req.user.alanyaID},
\code{parseInt(req.params.id, 10)}, \code{pool.execute}, codes 400/403/404/409, et les
utilitaires existants \code{sanitizeUrl} et \code{maskPresenceIfBlocked}
(\code{src/utils/blockUtils.js}).

\begin{center}
\begin{tabularx}{\linewidth}{@{}p{3.5cm}X@{}}
\toprule
\thd{Handler} & \thd{Rôle} \\
\midrule
\code{getLists}       & Listes du propriétaire + \code{member\_count}
                        (\code{LEFT JOIN \ldots{} GROUP BY}) \\
\code{createList}     & \code{\{ name, color \}} --- 409 si \code{uq\_list\_name} est violé \\
\code{updateList}     & Renommer / recolorer, avec contrôle de propriété \\
\code{deleteList}     & Le \code{CASCADE} supprime les membres \\
\code{getListMembers} & \code{JOIN} vers \code{users}, même projection que
                        \code{getPreferredContacts} \\
\code{addMember}      & Vérifie la propriété \textbf{et} que l'ami est un favori --- sinon 403 \\
\code{removeMember}   & --- \\
\bottomrule
\end{tabularx}
\end{center}

\section{Routes --- \code{src/routes/contactLists.js}}

Router Express avec \code{auth} sur \textbf{chaque} route et des blocs \code{@swagger},
comme \code{src/routes/contacts.js}.

\begin{center}
\begin{tabular}{@{}ll@{}}
\toprule
\mGET{}~\code{/}                          & \code{getLists} \\
\mPOST{}~\code{/}                         & \code{createList} \\
\mPUT{}~\code{/:idList}                   & \code{updateList} \\
\mDEL{}~\code{/:idList}                   & \code{deleteList} \\
\mGET{}~\code{/:idList/members}           & \code{getListMembers} \\
\mPOST{}~\code{/:idList/members/:friendID}& \code{addMember} \\
\mDEL{}~\code{/:idList/members/:friendID} & \code{removeMember} \\
\bottomrule
\end{tabular}
\end{center}

Montage dans \code{server.js}~: \code{app.use('/api/contact-lists', contactListRoutes)}.

\begin{constat}[Aucun endpoint de conversation à créer]
L'action « créer un groupe depuis une liste » réutilise
\mPOST{}~\code{/api/conversations/group}, déjà en service et déjà appelé côté mobile par
\code{createGroup} (\code{lib/api/chat\_api.dart}, ligne~27, signature
\code{\{participantIDs, groupName, groupPhoto\}}).
\end{constat}

\chapter{Client Flutter}

\section{Couche données}

\begin{center}
\begin{tabularx}{\linewidth}{@{}p{5.6cm}X@{}}
\toprule
\thd{Fichier} & \thd{Contenu} \\
\midrule
\code{lib/api/contact\_lists\_api.dart} \emph{(nouveau)} &
  Extension \code{ContactListsApi on TalkyApiClient}, \code{part of} le client, sur le
  modèle de \code{lib/api/users\_api.dart}. Déclarer le \code{part} dans
  \code{talky\_api\_client.dart}. \\
\code{lib/talky\_models.dart} &
  Modèle \code{ContactList} : \code{idList}, \code{name}, \code{color},
  \code{memberCount}, \code{fromJson} / \code{toJson}. \\
\code{lib/core/db/app\_database.dart} &
  Tables \code{LocalContactLists} et \code{LocalContactListMembers}
  (clé composite \code{\{idList, idFriend\}}). \\
\code{lib/core/services/local\_cache\_repository.dart} &
  Section « listes de contacts », calquée sur \code{watch/syncPreferredContacts}. \\
\bottomrule
\end{tabularx}
\end{center}

\begin{alerte}[Version du schéma local --- v15 devient v16]
La spécification d'origine annonçait un passage de 14 à 15. \code{schemaVersion} vaut
aujourd'hui \textbf{14} (\code{lib/core/db/app\_database.dart}, ligne~253), et le socle de
compte l'amène déjà à \textbf{15}. Les listes de contacts prennent donc \textbf{16}~:

\begin{lstlisting}[language=Dart, numbers=none]
if (from < 16) {
  await m.createTable(localContactLists);
  await m.createTable(localContactListMembers);
}
\end{lstlisting}

Puis régénérer~: \code{dart run build\_runner build -{}-delete-conflicting-outputs}.
La montée doit passer par \code{onUpgrade}, jamais par une recréation.
\end{alerte}

Le dépôt local expose~:
\code{watchContactLists()} (flux Drift, hors ligne d'abord),
\code{syncContactLists()} (même logique \code{previousIds} / \code{newIds} de
diff-suppression que \code{syncPreferredContacts}),
\code{watchListMembers(idList)} et \code{syncListMembers(idList)}.
Les mutations passent en ligne par l'API puis re-synchronisent, comme
\code{addContact} / \code{removeContact}. Les deux tables sont ajoutées à
\code{clearSession()}.

\section{Interface}

\begin{center}
\begin{tabularx}{\linewidth}{@{}p{6.2cm}X@{}}
\toprule
\thd{Fichier} & \thd{Objet} \\
\midrule
\code{screens/profile/preferred\_contacts\_screen.dart} &
  Rangée de puces au-dessus de la liste ; état \code{\_selectedListId}
  (\code{null} = Tous) combiné à la recherche existante
  (\code{filterUsersBySearch}). \\
\code{screens/profile/contact\_lists\_screen.dart} \emph{(nouveau)} &
  CRUD des listes. \\
\code{screens/profile/contact\_list\_detail\_screen.dart} \emph{(nouveau)} &
  Membres, ajout / retrait, bouton « Créer un groupe ». \\
\code{screens/profile/profile\_screen.dart} &
  Entrée « Listes de contacts ». \\
\code{lib/l10n/*} &
  \code{contactLists}, \code{createList}, \code{listName}, \code{renameList},
  \code{deleteList}, \code{addToList}, \code{removeFromList},
  \code{createGroupFromList}, \code{noLists}\ldots{} \\
\bottomrule
\end{tabularx}
\end{center}

\begin{note}[Composants à réutiliser, pas à réécrire]
Le patron de puce existe déjà~: \code{\_buildFilterChip} dans
\code{lib/screens/chats/chats\_screen.dart}, \textbf{ligne~584} --- et non 539--569 comme
l'indiquait la spécification d'origine, dont la référence avait vieilli.
Pour le reste~: \code{AppAvatar}, \code{AppSearchField}, \code{EmptyState},
\code{AppBottomSheet}, et les jetons \code{AppSpacing} / \code{AppRadius} /
\code{context.colors} / \code{context.l10n}.
\end{note}

\chapter{Arbitrage : liste de contacts \emph{n'est pas} audience de diffusion}

La conception initiale (\code{alanya-contexte-conception.md}, §3) proposait \emph{une seule}
abstraction d'audience, réutilisée pour les segments de messages officiels, les segments de
statuts officiels et les listes de contacts préférés --- « ne pas la dupliquer ».

\textbf{Décision prise~: deux modèles distincts.} Les deux objets n'ont ni le même
propriétaire, ni la même échelle, ni la même nature.

\begin{center}
\begin{tabularx}{\linewidth}{@{}p{3.1cm}XX@{}}
\toprule
 & \thd{Liste de contacts} & \thd{Audience de diffusion} \\
\midrule
Propriétaire & un utilisateur & un administrateur \\
Membres      & ses favoris, énumérés & tous les comptes, décrits par un critère \\
Échelle      & quelques dizaines & potentiellement toute la base \\
Nature       & des lignes & une \emph{requête} \\
Usage        & filtrer, créer un groupe & cibler une diffusion \\
\bottomrule
\end{tabularx}
\end{center}

Les fusionner obligerait à choisir entre deux mauvaises options~:

\begin{enumerate}[leftmargin=1.6em]
  \item \textbf{matérialiser} l'audience de diffusion en lignes d'appartenance --- c'est
        exactement le \emph{fan-out on write} que le §9 du même document rejette~;
  \item \textbf{doter} les listes personnelles d'un moteur de critères dont elles n'auront
        jamais l'usage.
\end{enumerate}

\begin{note}[Ce qui reste mutualisé]
Deux schémas distincts, mais \textbf{un seul résolveur} côté serveur~: une fonction qui
prend l'un ou l'autre et renvoie la liste des destinataires. La mutualisation se fait au
niveau du code d'exécution, où elle est utile, pas au niveau du schéma, où elle coûte.
\end{note}

\chapter{Vérification}

\section{Backend}

Appliquer \code{024\_contact\_lists.sql} sur MySQL, redémarrer, puis tester chaque endpoint
via Swagger ou \code{curl} avec un jeton valide~:

\begin{enumerate}[leftmargin=1.6em]
  \item créer une liste $\rightarrow$ 201~;
  \item recréer le même nom $\rightarrow$ 409~;
  \item ajouter un contact \textbf{non favori} $\rightarrow$ 403~;
  \item lister les membres, en retirer un, supprimer la liste~;
  \item vérifier que la suppression de la liste n'a supprimé \textbf{aucun} favori.
\end{enumerate}

\section{Client}

\code{flutter pub get} $\rightarrow$
\code{dart run build\_runner build -{}-delete-conflicting-outputs} $\rightarrow$
\code{flutter run}.

\begin{enumerate}[leftmargin=1.6em]
  \item Créer « Famille », y ajouter trois favoris~; la puce « Famille » filtre la liste.
  \item « Créer un groupe » ouvre une conversation de groupe avec ces membres et le nom
        pré-rempli.
  \item \textbf{Migration v15 $\rightarrow$ v16} sur un appareil déjà installé~: la mise à
        jour ne doit \textbf{pas} vider les données existantes.
  \item \textbf{Hors ligne}~: réseau coupé, listes et membres restent visibles depuis le
        cache Drift~; les mutations affichent un état hors ligne cohérent.
  \item Retirer un favori appartenant à une liste, et vérifier le comportement retenu.
\end{enumerate}

\section{Points ouverts}

\begin{enumerate}[leftmargin=1.6em]
  \item Retrait d'un favori encore présent dans des listes~: cascade applicative, ou filtrage
        à la lecture~?
  \item Une limite au nombre de listes par utilisateur~? Rien ne l'impose techniquement,
        mais la rangée de puces se dégrade au-delà d'une dizaine.
  \item Extensions volontairement hors périmètre, possibles sans changer les tables~:
        appel de groupe depuis une liste, diffusion de statut restreinte à une liste.
\end{enumerate}
